Three limits are worth stating plainly. First, the study is simulation-only. We argue the
offline-to-online mapping is the very thing a real-robot user needs before committing to
deployment evaluation, and sim-to-real correlation work supports that conclusions transfer in
direction, but we do not run hardware. Second, one architecture. Diffusion Policy is the
modal manipulation policy, yet the gap could behave differently for autoregressive or
vision-language-action models; an ACT replication fits in one extra session and is the
obvious next step. Third, statistical power. With three seeds and a Bonferroni-corrected
$\alpha = 0.0125$, the protocol has roughly \PowerHone{} power to detect a ten-percentage-point
gap and is well-powered only for larger effects; we report the realised power curve rather
than a single optimistic figure, and we treat the weakest-powered test, H3, with
corresponding caution. None of these undercuts the core deliverable: a locked, reproducible
framework for asking whether an offline metric can be trusted.
